% !TEX program = xelatex
% ============================================================
%  واجب منزلي — Homework Assignment Template
%  قالب واجب منزلي للطلاب مع ترويسة بمعلومات المقرر والطالب
%  وأسئلة مرقّمة مع مساحة للإجابة (رياضيات، إجابة قصيرة، مقال).
%  Compile with: xelatex main.tex (run twice)
% ============================================================
%  Features:
%    - A4, 11pt, article class
%    - Header block: course, student, assignment #, due date
%    - Numbered problems via enumerate with answer space
%    - Math problem, short-answer problem, essay problem examples
%    - Lined answer area and boxed answer area helpers
%    - fancyhdr header/footer
% ============================================================

\documentclass[11pt,a4paper]{article}

% ---- Geometry ----
\usepackage[a4paper,margin=2.5cm,top=2.5cm,bottom=2.5cm,headheight=16pt]{geometry}

% ---- Core packages ----
\usepackage{graphicx}
\usepackage{array}
\usepackage{booktabs}
\usepackage{tabularx}
\usepackage{multirow}
\usepackage{enumitem}
\usepackage{float}
\usepackage{titlesec}
\usepackage{fancyhdr}
\usepackage{setspace}
\usepackage{xcolor}
\usepackage{amsmath}

% ---- RTL / Arabic language support (order matters) ----
\usepackage{bidi}
\usepackage[no-math]{fontspec}
\usepackage[quiet]{polyglossia}

% ---- Fonts ----
\setmainfont[Scale=1]{Amiri-Regular.ttf}
\newfontfamily\arabicfont[Script=Arabic,Scale=0.95]{Amiri-Regular.ttf}
\newfontfamily\englishfont[Scale=0.95]{TeX Gyre Termes}

% ---- Languages ----
\setdefaultlanguage[calendar=gregorian,hijricorrection=1]{arabic}
\setotherlanguage[variant=british]{english}

% ---- Hyperref (load last, after polyglossia/bidi) ----
\usepackage[breaklinks,hidelinks]{hyperref}

% ---- Colors ----
\definecolor{hdraccent}{HTML}{1F3A5F}
\definecolor{boxline}{HTML}{888888}

% ---- Line spacing ----
\setstretch{1.4}

% ---- Captions in Arabic ----
\addto\captionsarabic{%
  \renewcommand{\figurename}{الشكل}%
  \renewcommand{\tablename}{الجدول}%
}

% ---- Section formatting ----
\titleformat{\section}{\Large\bfseries}{\thesection}{0.7em}{}
\titleformat{\subsection}{\large\bfseries}{\thesubsection}{0.6em}{}
\titlespacing*{\section}{0pt}{1.2ex plus 0.3ex}{0.8ex plus 0.2ex}

% ---- Headers / footers ----
\pagestyle{fancy}
\fancyhf{}
\renewcommand{\headrulewidth}{0.4pt}
\renewcommand{\footrulewidth}{0pt}
% TODO: عدّل اسم المقرر ورقم الواجب في الترويسة.
\fancyhead[R]{\small واجب منزلي (Homework)}
\fancyhead[L]{\small اسم المؤسسة (Institution Name)}
\fancyfoot[C]{\small \thepage}

% ============================================================
% Answer-space helpers
% ============================================================

% \answerlines{height} — lined answer area (height in cm)
\newcommand{\answerlines}[1]{%
  \par\vspace{0.3em}%
  \dimen0=#1\relax
  \loop
    \noindent\hrulefill\par\vspace{0.7em}%
    \advance\dimen0 by -1.1em\relax
  \ifdim\dimen0>1.1em\repeat
  \par
}

% \answerbox{height} — empty boxed answer area
\newcommand{\answerbox}[1]{%
  \par\vspace{0.3em}%
  \noindent\fbox{\rule{0pt}{#1}\rule{\dimexpr\linewidth-2\fboxsep-2\fboxrule\relax}{0pt}}%
  \par
}

\begin{document}

% ============================================================
% TITLE / HEADER BLOCK
% ============================================================
\begin{center}
{\LARGE\bfseries واجب منزلي (Homework Assignment)}\\[0.4em]
{\large\bfseries اسم المؤسسة (Institution Name)}
\end{center}

\vspace{0.3em}
{\color{hdraccent}\rule{\linewidth}{1.2pt}}
\vspace{0.5em}

% ---- Info table ----
\noindent
\begin{tabularx}{\linewidth}{@{}p{0.30\linewidth}X@{}}
\textbf{المقرر (Course):} & % TODO: ضع اسم المقرر هنا
اسم المقرر (Course Name) \\
\textbf{اسم الطالب (Student):} & % TODO: ضع اسم الطالب هنا
\hrulefill \\
\textbf{الرقم الجامعي (ID):} & % TODO: ضع الرقم الجامعي هنا
\hrulefill \\
\textbf{رقم الواجب (Assignment \#):} & % TODO: ضع رقم الواجب هنا
\hrulefill \\
\textbf{تاريخ التسليم (Due Date):} & % TODO: ضع تاريخ التسليم هنا
\hrulefill \\
\textbf{اسم المعلم (Instructor):} & % TODO: ضع اسم المعلم هنا
\hrulefill \\
\end{tabularx}

\vspace{0.4em}
{\color{hdraccent}\rule{\linewidth}{0.8pt}}
\vspace{0.8em}

% ============================================================
% PROBLEMS
% ============================================================
\section*{الأسئلة (Problems)}
% TODO: استبدل الأسئلة أدناه بأسئلة الواجب الفعلية.

\begin{enumerate}[label=\textbf{\arabic*.}, leftmargin=2.2em, itemsep=1.2em]

% ---- Problem 1: Math ----
\item \textbf{مسألة رياضيات (Math Problem):}\\
% TODO: ضع المسألة الرياضية هنا.
احسب قيمة $x$ في المعادلة التالية ثم اعرض خطوات الحل كاملة:
\[
3x + 7 = 2x + 22
\]
\textit{الحل (Solution):}
\answerlines{4cm}

% ---- Problem 2: Short answer ----
\item \textbf{إجابة قصيرة (Short-Answer Problem):}\\
% TODO: ضع سؤال الإجابة القصيرة هنا.
اذكر ثلاثة من العوامل المؤثرة في سرعة التفاعل الكيميائي مع شرح موجز لكل عامل.
\answerlines{4cm}

% ---- Problem 3: Essay ----
\item \textbf{سؤال مقالي (Essay Problem):}\\
% TODO: ضع سؤال المقال هنا.
ناقش أثر الثورة الصناعية على النمو الاقتصادي والاجتماع في أوروبا، موضحاً
أهم التغيّرات التي طرأت على أنماط الحياة والعمل. اكتب إجابتك في الإطار أدناه.
\answerbox{6cm}

\end{enumerate}

\vspace{1em}
{\color{hdraccent}\rule{\linewidth}{0.6pt}}
\vspace{0.3em}
\begin{center}
\textit{بالتوفيق للجميع (Good luck)}
\end{center}

\end{document}
